\begin{ex}
Tính giá trị của biểu thức: $42 + 35 - 24$.
\shortans{53}
\loigiai{
Thực hiện: $42 + 35 = 77$; $77 - 24 = 53$.
\textbf{Đáp số:} $53$.
}
\end{ex}

\begin{ex}
Tính giá trị của biểu thức: $88 - 53 + 14$.
\shortans{49}
\loigiai{
Thực hiện: $88 - 53 = 35$; $35 + 14 = 49$.
\textbf{Đáp số:} $49$.
}
\end{ex}

\begin{ex}
Tính giá trị của biểu thức: $99 - 44 - 33$.
\shortans{22}
\loigiai{
Thực hiện: $99 - 44 = 55$; $55 - 33 = 22$.
\textbf{Đáp số:} $22$.
}
\end{ex}

\begin{ex}
Tìm số tự nhiên lớn nhất có hai chữ số.
\shortans{99}
\loigiai{
Số tự nhiên lớn nhất có hai chữ số là số $99$.
\textbf{Đáp số:} $99$.
}
\end{ex}

\begin{ex}
Tìm số tự nhiên bé nhất có hai chữ số khác nhau.
\shortans{10}
\loigiai{
Số bé nhất có hai chữ số khác nhau là số $10$.
\textbf{Đáp số:} $10$.
}
\end{ex}

\begin{ex}
Tìm số liền sau của số 99.
\shortans{100}
\loigiai{
Số liền sau của 99 là số $100$.
\textbf{Đáp số:} $100$.
}
\end{ex}

\begin{ex}
Tìm số liền trước của số 80.
\shortans{79}
\loigiai{
Số liền trước của 80 là số $79$.
\textbf{Đáp số:} $79$.
}
\end{ex}

\begin{ex}
Điền số thích hợp: $54 + \dots = 89$.
\shortans{35}
\loigiai{
Số cần điền là: $89 - 54 = 35$.
\textbf{Đáp số:} $35$.
}
\end{ex}

\begin{ex}
Điền số thích hợp: $76 - \dots = 42$.
\shortans{34}
\loigiai{
Số cần điền là: $76 - 42 = 34$.
\textbf{Đáp số:} $34$.
}
\end{ex}

\begin{ex}
Một trường có 55 học sinh nam và 43 học sinh nữ khối 1. Hỏi khối 1 có tất cả bao nhiêu học sinh?
\shortans{98}
\loigiai{
Tổng số học sinh khối 1 là: $55 + 43 = 98$ (học sinh).
\textbf{Đáp số:} $98$.
}
\end{ex}

\begin{ex}
Số học sinh nam nhiều hơn số học sinh nữ ở bài toán trên là bao nhiêu bạn?
\shortans{12}
\loigiai{
Số học sinh nam nhiều hơn nữ là: $55 - 43 = 12$ (học sinh).
\textbf{Đáp số:} $12$.
}
\end{ex}

\begin{ex}
Một cửa hàng có 85 chiếc xe đạp, đã bán được 53 chiếc xe đạp. Hỏi cửa hàng còn lại bao nhiêu chiếc xe đạp?
\shortans{32}
\loigiai{
Số xe đạp còn lại là: $85 - 53 = 32$ (chiếc xe đạp).
\textbf{Đáp số:} $32$.
}
\end{ex}

\begin{ex}
Một sợi dây dài 95 cm. Người ta cắt đi hai đoạn dài 30 cm và 25 cm. Sợi dây còn lại dài bao nhiêu xăng-ti-mét?
\shortans{40}
\loigiai{
Tổng độ dài đã cắt: $30 + 25 = 55$ (cm).\nĐộ dài còn lại là: $95 - 55 = 40$ (cm).
\textbf{Đáp số:} $40$.
}
\end{ex}

\begin{ex}
Đồng hồ chỉ 10 giờ đêm (tức 22 giờ). Sau 2 giờ nữa là mấy giờ đêm (ngày mới)?
\shortans{12}
\loigiai{
Lúc đó là: $10 + 2 = 12$ (giờ đêm, tức 0 giờ ngày mới).
\textbf{Đáp số:} $12$.
}
\end{ex}

\begin{ex}
Nếu hôm nay là Thứ Năm ngày 14 thì Thứ Năm tuần trước là ngày mùng mấy?
\shortans{7}
\loigiai{
Thứ Năm tuần trước là ngày: $14 - 7 = 7$.
\textbf{Đáp số:} $7$.
}
\end{ex}

\begin{ex}
Nếu hôm nay là Thứ Bảy ngày 21 thì Thứ Bảy tuần sau là ngày bao nhiêu?
\shortans{28}
\loigiai{
Thứ Bảy tuần sau là ngày: $21 + 7 = 28$.
\textbf{Đáp số:} $28$.
}
\end{ex}

\begin{ex}
Một hình chữ nhật có chiều dài 30 cm, chiều rộng 20 cm. Chiều dài hơn chiều rộng bao nhiêu xăng-ti-mét?
\shortans{10}
\loigiai{
Chiều dài hơn chiều rộng là: $30 - 20 = 10$ (cm).
\textbf{Đáp số:} $10$.
}
\end{ex}

\begin{ex}
Khối lập phương có bao nhiêu mặt phẳng?
\shortans{6}
\loigiai{
Khối lập phương có đúng 6 mặt phẳng.
\textbf{Đáp số:} $6$.
}
\end{ex}

\begin{ex}
Khối hộp chữ nhật có bao nhiêu cạnh?
\shortans{12}
\loigiai{
Khối hộp chữ nhật có đúng 12 cạnh.
\textbf{Đáp số:} $12$.
}
\end{ex}

\begin{ex}
Tìm hai số có tổng bằng 77 và hiệu bằng 11. Số lớn hơn là số mấy?
\shortans{44}
\loigiai{
Hai số là 44 và 33 ($44 + 33 = 77$ và $44 - 33 = 11$). Số lớn hơn là $44$.
\textbf{Đáp số:} $44$.
}
\end{ex}

\begin{ex}
Tìm hai số có tổng bằng 88 và hiệu bằng 22. Số bé hơn là số mấy?
\shortans{33}
\loigiai{
Hai số là 55 và 33 ($55 + 33 = 88$ và $55 - 33 = 22$). Số bé hơn là $33$.
\textbf{Đáp số:} $33$.
}
\end{ex}

\begin{ex}
Có bao nhiêu số tròn chục có hai chữ số?
\shortans{9}
\loigiai{
Có 9 số tròn chục: $10, 20, 30, 40, 50, 60, 70, 80, 90$.
\textbf{Đáp số:} $9$.
}
\end{ex}

\begin{ex}
Có bao nhiêu số có hai chữ số giống nhau?
\shortans{9}
\loigiai{
Có 9 số: $11, 22, 33, 44, 55, 66, 77, 88, 99$.
\textbf{Đáp số:} $9$.
}
\end{ex}

\begin{ex}
Một cuộn ruy băng dài 90 cm. Bạn Linh cắt lấy 3 đoạn, mỗi đoạn dài 20 cm. Hỏi cuộn ruy băng còn lại dài bao nhiêu xăng-ti-mét?
\shortans{30}
\loigiai{
Tổng độ dài 3 đoạn cắt là: $20 + 20 + 20 = 60$ (cm).\nĐộ dài còn lại là: $90 - 60 = 30$ (cm).
\textbf{Đáp số:} $30$.
}
\end{ex}

\begin{ex}
Tính giá trị của biểu thức: $99 - 11 - 22 - 33$.
\shortans{33}
\loigiai{
Thực hiện: $99 - 11 = 88$; $88 - 22 = 66$; $66 - 33 = 33$.
\textbf{Đáp số:} $33$.
}
\end{ex}
